\section{Related work and comparative scope}
\label{sec:related}
\paragraph{Interposing beneath an operating system.}
SubVirt demonstrates placing a VMM below existing Windows and Linux systems
in a security-threat setting~\cite{subvirt}. Blue Pill describes using AMD SVM
to move a running OS beneath a thin monitor without rebooting~\cite{bluepill}.
These precedents rule out claiming reboot-free OS interposition itself as our
invention. Our authorized laboratory use also makes no claim of undetectability.
The engineering question here is the interaction between such interposition,
an already-present outer Hyper-V, a later unmodified intermediate VMM, and
explicit page-resource retirement.

\paragraph{Nested execution and lower-layer control.}
Turtles virtualizes VMX and composes translation layers to execute unmodified
guest hypervisors~\cite{turtles}. \ks{} uses this established model for virtual
VMCS state, exit reflection, and composed EPT. CloudVisor places a small monitor
below a commodity VMM and separates VM protection from resource
management~\cite{cloudvisor}. Our page control shares the broad placement of
authority below the intermediate VMM, but not CloudVisor's protection goal.
\ks{} trusts the outer platform and architectural cooperation from the
intermediate VMM, requires privileged guest page preparation in the experiment,
and supplies no confidentiality/integrity guarantee against a compromised VMM.

\paragraph{Replacing a running hypervisor.}
HyperFresh uses a pre-existing hyperplexor and memory remapping to transfer
running VMs to a replacement hypervisor on the same machine~\cite{hyperfresh}.
HyperTP combines migration-based replacement and in-place microreboot using a
unified representation of VM state, with Xen/KVM implementations~\cite{hypertp}.
These systems transfer already-running VM state. Our measured operation instead
inserts a monitor with descendant VMMs absent and starts VMware later. We do not
equate those operations or compare their published downtime numbers with our
Windows-only driver-body times.

\paragraph{Nested performance.}
HyperTurtle moves selected nested critical paths into eBPF programs executed by
the bare-metal hypervisor to reduce world switches~\cite{hyperturtle}. Our
prototype does not install such an outer-hypervisor extension. Its cost profile
motivates examining boundary transitions beyond the C dispatcher, but our
experiments do not isolate all of those costs and do not measure a performance
advantage over HyperTurtle. No prior system is treated as a benchmark baseline
without a comparable local run.

\paragraph{Application value.}
The HTTP comparison establishes a concrete distinction in authority and recovery
state: after cooperative page setup, a containing-Windows operator can serve
private backing while preserving the original page, whereas the demonstrated
direct-write alternative requires guest write permission and saved contents.
This can support controlled fault campaigns in nested test environments that
use an unmodified VMM. Whether the additional monitor is worthwhile depends on
those control requirements and its measured overhead. Broader application value
requires workloads with meaningful recovery behavior and comparisons against
other available fault-injection interfaces.

\section{Limitations and threats to validity}
\label{sec:limitations}
\paragraph{Insertion scope.}
The measured insertion runs with descendant VMMs absent, and VMware starts
afterward. We do not move an already-running VMware/TinyCore stack beneath the
monitor, and we report no measurement of such a transfer. In the inner Hyper-V
topology, admission is refused at capability discovery rather than at
\code{VMXON}: with the monitor not resident, the inner Windows root partition
reports no \code{CPUID.1:ECX[5]} VMX bit and no nested capability MSRs, and the
driver declines the request before submission. Because the monitor is not
resident during that reading, the absent VMX bit is not self-inflicted. The
refusal is the inner hypervisor's decision for its own root partition in this
topology and build; we do not generalize it to other topologies, and we do not
attribute it to a vendor defect.

\paragraph{Page semantics.}
Process ownership and sampled source-path identity improve lifetime accounting
but cannot detect identical-bit ABA, all guest reboots, snapshot restoration,
or application-object reuse at the same GPA. The caller must manage those
boundaries. A region admitted by scanning is weaker still, and deliberately so:
its recheck visits one page per 4096 shadow fills, so a source entry that stops
agreeing is detected within a bounded number of samples rather than at the
composition that follows it. Callers who cannot tolerate that interval should
use the coarse-source rule, which needs no recheck because a single source leaf
cannot disagree with itself.

A whole-page readback on two CPUs is not proof of atomic replacement
for arbitrary concurrent readers/writers, DMA, MMIO, or executable-page
workloads. Successful synchronous retirement and sampled readback are different
forms of evidence; neither should be enlarged into universal stale-mapping
freedom without the stated assumptions.

\paragraph{A build artifact that presented as a descendant hang.}
Exercising regions produced a failure worth reporting for how it looked rather
than for what it was. A descendant under a memory workload stopped executing:
its VMM process kept serving its framebuffer and answering the host, consumed no
CPU, ignored console input, and could not be stopped by its own command-line
tool. From outside that is indistinguishable from a hang. It was not one. The
VMM's log names it exactly --- \code{MONITOR PANIC: VM-entry failed; VMCS valid
(error code 7)}, an entry rejected for invalid control fields --- and the
unrecoverable-error dialog that would have said so is invisible for a VM started
from session~0.

The cause is not in the sources. Two driver binaries built from identical
sources differ only in how they were built: \code{msbuild /t:Build} produced one
that failed this way in three runs of three, and \code{/t:Rebuild} produced one
that ran an 800\,MiB allocation, a snapshot cycle and a further allocation with
none. We do not explain why an incremental build differs. We record that it did,
that the two binaries have different hashes, and that a reading taken with a
driver that was not cleanly rebuilt establishes nothing.

Two observations from the same session remain unexplained and are not folded
into that one. The containing Windows VM itself once lost its heartbeat and
ignored an injected non-maskable interrupt, producing no bugcheck and therefore
no dump --- consistent with this build installing a private host
interrupt-descriptor table, so the usual convert-a-hang-into-a-dump step is
unavailable here. And after a control process was killed while a removal was in
flight, three retried removals returned \code{STATUS\_HV\_OPERATION\_FAILED}
with \code{retired=1}, so the backing stayed retained: the designed fail-closed
path behaving as specified, but never subsequently succeeding. The kill is a
plausible cause we cannot exclude, and a later removal after an ordinary
revocation drained cleanly on the first attempt.

\paragraph{Coverage and reliability.}
Current page-cycle counts are small, and software-injected allocation/drain
failures do not emulate all hardware failures. Source-path comparison has
portable adversarial coverage but lacks deliberate hardware EPT mutation tests.
The region lease has been revoked by a real source change on hardware: a
snapshot of the descendant taken while a 2\,MiB region was published drops the
intermediate VMM's entries for that region entirely, and the lease was revoked
within five seconds and removed cleanly. It reported translation-changed rather
than region-drifted, and that separation is structural rather than incidental:
the always-on check covers the path the lease was captured on and runs at every
shadow fill, so it sees any whole-table event first, while the region recheck
exists for the case it cannot see --- one page of the region losing permission
while the captured path stays identical. Nothing this VMM can be asked to do
produces that case, so the region-drifted path itself remains covered only by
the portable suite and its eight injected defects.
High-load, write-isolation, and owner-exit results belong to the historical
build. A historical two-by-two configuration has a 599.43-second observation
with 21 samples; there is no new ten-minute or overnight soak for the current
four-by-two build. That older duration cannot establish current crash freedom,
bounded long-term memory growth, or whole-driver leak freedom.

\paragraph{Performance generality.}
One heterogeneous physical CPU, one outer Hyper-V build, and one current
Windows~1 build provide no cross-machine estimate. Guest CPU affinity does not
eliminate physical scheduling or host activity. The current performance suite
has five measured repetitions and no descendant VMMs; the broader historical
suite has block-order drift. The optimized nested comparison crosses driver
versions and boots. Persistent descendant disk performance, external networking,
and active-lease steady-state overhead remain unmeasured. A matched no-monitor
VMware backend is a different case, and we separate it rather than folding it
into that list: two attempts are recorded with the driver service stopped and no
monitor resident. One started a \code{vmware-vmx} process and wrote a new VMware
log but left guest boot unverified; the other was refused at start with a
cancellation, leaving no new process and no new log. Each attempt retains the
Windows boot identity it ran under. These records establish that the comparison
was attempted and did not yield a usable baseline; they do not attribute either
outcome to the monitor's absence, and no no-monitor descendant comparison is
reported. Windows 10 LTSC, other CPU families, and other Windows builds
have not been evaluated in the current matrix.

\paragraph{Application observation.}
The HTTP oracle validates a controlled fixed schema, not production business
logic. Its current dataset lacks an independent holder-liveness probe at every
response because the measured collector emitted a constant field. The raw data
and this defect remain visible; fresh PFN observations support only the narrower
claim stated in Section~\ref{sec:application}. A bounded observer expiration is
an incomplete trial. Settling exclusions in the direct-write comparator prevent
an atomicity or matched-latency interpretation. Paced closed-loop Windows TCP
observations likewise do not measure saturated throughput or an exact pause.

\paragraph{Artifact scope.}
Per-run logs, source/binary identities, analysis scripts, and generated tables
make claims inspectable. They do not independently reproduce the original
hardware runs. The separately prepared redacted archive contains evidence and
analysis, not a complete anonymous buildable system image. Removing direct
identity strings also does not prevent linking a public project to its hardware,
dates, or distinctive measurements.

\section{Conclusion}
\ks{} demonstrates a bounded way to introduce a VMX monitor beneath a running
Hyper-V Windows guest and subsequently control selected memory in a VMware
descendant. Its page lease separates admission, publication, observed revocation,
CPU invalidation, and reclamation, and its experiments connect those events to
multicore readback and an HTTP policy effect. The evaluated pages can be
substituted and restored with tracked resources reclaimed, and so can a 2\,MiB
region of them, admitted on a property that makes one leaf safe for all of it
rather than on a proxy for that property.

The limits are as concrete as the capability. Significant exit and RTT costs
remain; every measured insertion begins with descendant VMMs absent; the inner
Hyper-V root partition refuses admission in the tested configuration; a 1\,GiB
leaf is refused on this stack, because the intermediate VMM maps at 4\,KiB and a
region that size exceeds a bound this design keeps deliberately; which bases can
be replaced at all follows what the descendant has touched rather than what it
was configured with; the region recheck's own revocation reason has never fired
on hardware, because the always-on check sees every source change this VMM can
be made to produce; and two failures found while exercising regions remain
unexplained, after a third turned out to be a driver built incrementally rather
than cleanly. These define the next work: existing-state handoff, better
attribution of boundary costs, stronger page-identity and concurrency
validation, a source change confined to one page of a region, an explanation for
those two failures, and independent hardware and application evaluation.

\section*{Data, code, and review status}
The implementation and versioned evidence are in the public \ks{} repository;
Appendix~\ref{app:artifact} identifies exact revisions and reproduction commands.
Tables in this preprint are generated from archived summaries with recorded
input SHA256 values. \path{docs/next/arxiv/build-validation.json} records, for
the exact build it names, which checks were performed and which were not; the
author's review of this revision is recorded there as pending, and earlier
revisions were reviewed and their data verified. This preprint has not been peer
reviewed or submitted to arXiv as part of its preparation, and it does not
assert conference acceptance or a publication DOI.

\section*{Use of AI assistance}
AI coding assistants -- OpenAI Codex, and Anthropic Claude for the region
mechanism, its portable tests and the corresponding manuscript sections --
assisted with prototype implementation and debugging, experiment and analysis
scripts, authorized laboratory execution, evidence archiving, and drafting the
manuscript and diagrams. Reported measurements are tied to retained
logs; unmeasured behavior is identified separately. The author has reviewed the
source, experimental procedures, raw data, statistical interpretation, citations,
and claims for this draft. The AI system is not an author.
